\documentclass[11pt]{article}
\usepackage[margin=1in]{geometry}
\usepackage{amsmath,amssymb}
\usepackage{xcolor}
\usepackage{booktabs}
\usepackage{enumitem}
\usepackage{fancyhdr}
\usepackage{tcolorbox}
\tcbuselibrary{skins,breakable}
\usepackage{hyperref}
\hypersetup{colorlinks=true, linkcolor=blue, urlcolor=blue}

\definecolor{primary}{RGB}{30,64,140}
\definecolor{tipbg}{RGB}{235,244,255}
\definecolor{examplebg}{RGB}{240,248,235}
\definecolor{formulabg}{RGB}{255,247,224}
\definecolor{warnbg}{RGB}{255,236,236}

\pagestyle{fancy}
\fancyhf{}
\renewcommand{\headrulewidth}{0.4pt}
\cfoot{\thepage}

\newtcolorbox{tipbox}[1][Tip]{
  colback=tipbg, colframe=primary, fonttitle=\bfseries,
  title=#1, breakable, rounded corners
}
\newtcolorbox{examplebox}[1][Example]{
  colback=examplebg, colframe=primary!70!black, fonttitle=\bfseries,
  title=#1, breakable, rounded corners
}
\newtcolorbox{formulabox}[1][Formula]{
  colback=formulabg, colframe=primary!70!black, fonttitle=\bfseries,
  title=#1, breakable, rounded corners
}
\newtcolorbox{warnbox}[1][Warning]{
  colback=warnbg, colframe=red!60!black, fonttitle=\bfseries,
  title=#1, breakable, rounded corners
}


\title{\color{primary}\textbf{Calendar --- Remaining Variants}\\[4pt]\large The Last \~{}5\%: Reasoning-Style Phrasings + a One-Page Decision Tree}
\author{Aptitude Grind}
\date{}

\lhead{\small Calendar --- Remaining Variants}
\rhead{\small Aptitude Grind}

\begin{document}
\maketitle
\thispagestyle{fancy}

\noindent Your first three booklets (\texttt{01-fundamentals}, \texttt{02-fast-tricks-and-shortcuts},
\texttt{03-question-pattern-catalog}) cover essentially all high-frequency TCS NQT Calendar
questions, and most of general placement Calendar too. This file adds the small remainder: ten
reasoning-style variants that show up occasionally (more often in Infosys/Wipro-style papers than
pure TCS), five extra trigger tricks, and a decision tree to pick the right method under time
pressure. Everything below still reduces to the core ideas in \texttt{01-fundamentals.tex} and the
Anchor-Day Method in \texttt{02-fast-tricks-and-shortcuts.tex}.

\section*{Category I --- Reasoning-Style Variants}

\subsection*{I1. Age + Calendar}
\begin{tipbox}[Phrasing]
``Ram was born on a Monday. His 18th birthday falls on a Wednesday. How many leap years occurred in
between (assuming no 29 February birthday)?''
\end{tipbox}
\textbf{Method:} Total weekday shift over 18 years = Wednesday $-$ Monday = 2 days forward. Each
ordinary year contributes 1 odd day, each leap year contributes 2. If $L$ is the number of leap
years, then $(18-L)(1)+L(2)\equiv2\pmod7 \Rightarrow 18+L\equiv2\pmod7 \Rightarrow L\equiv5\pmod7$.
Since $L$ must be a small number of leap years in an 18-year span (4 or 5 typically), \textbf{L = 5}.

\subsection*{I2. Date Mirror Questions}
\begin{tipbox}[Phrasing]
``If 12 March is a Tuesday, what day is 21 March?''
\end{tipbox}
\textbf{Method:} Never recount from the start of the month --- just take the difference in dates and
reduce mod 7. $21-12=9$ days; $9 \bmod 7=2$; Tuesday + 2 = \textbf{Thursday}. This is Trick 7 below,
applied directly.

\subsection*{I3. Consecutive Leap-Year Logic}
\begin{tipbox}[Phrasing]
``Three leap years occur between 1 January of year A and 1 January of year B (A earlier than B). If 1
January of year A is a Sunday, what day is 1 January of year B, given there are 10 years between A
and B?''
\end{tipbox}
\textbf{Method:} Same odd-day accumulation as I1, just phrased as a given leap-year count instead of
an asked one. Odd days $=(10-3)(1)+3(2)=7+6=13\equiv6\pmod7$; Sunday + 6 = \textbf{Saturday}. Always
check whether the leap-year count is \emph{given} (as here) or \emph{asked} (as in I1) --- the same
equation is just solved for a different variable.

\subsection*{I4. Financial / Office Calendar (Working-Day Count)}
\begin{tipbox}[Phrasing]
``An office is closed on all Sundays and on the second Saturday of every month. How many working days
are there in a 30-day month that starts on a Saturday?''
\end{tipbox}
\textbf{Method:} List the calendar out mod 7 for the Sundays and identify which Saturday is the
``second'' one; subtract both sets from the total days, watching for any date that is \emph{both} a
Sunday and a second-Saturday-adjacent date. With the month starting on Saturday: Saturdays fall on 1,
8, 15, 22, 29 (2nd Saturday = 8th); Sundays fall on 2, 9, 16, 23, 30 (5 Sundays). Total off-days
$=5\text{ (Sundays)}+1\text{ (second Saturday)}=6$; \textbf{working days = 30 $-$ 6 = 24}.

\subsection*{I5. Number of Sundays/Mondays Between Two Dates}
\begin{tipbox}[Phrasing]
``How many Sundays are there between 15 January and 15 April (inclusive of the range) of the same
year, given 15 January is a Thursday?''
\end{tipbox}
\textbf{Method:} Find the day-of-week of the range's start, locate the first Sunday, then count
forward in steps of 7 to the end. 15 Jan is Thursday $\Rightarrow$ first Sunday is 18 Jan. Days from
18 Jan to 15 Apr $=13\text{ (rest of Jan)}+28\text{ (Feb)}+31\text{ (Mar)}+15\text{ (Apr)}=87$ days
$\Rightarrow 87/7=12$ full weeks plus a remainder, giving \textbf{13 Sundays}. Always find the
\emph{first} occurrence of the target weekday in the range before dividing by 7 --- don't divide the
whole range by 7 directly.

\subsection*{I6. First/Last Working Day}
\begin{tipbox}[Phrasing]
``Find the first Monday after 15 June, given 15 June is a Friday.''
\end{tipbox}
\textbf{Method:} Find the weekday gap to the target day, then add it as a date offset. Friday to
Monday is $3$ days forward $\Rightarrow$ \textbf{18 June}. For a ``last weekday before a date''
version, the same idea works backward: find the gap going back to the most recent occurrence of the
target weekday.

\subsection*{I7. Month Identification from Weekday Counts}
\begin{tipbox}[Phrasing]
``A particular month has 5 Sundays, 5 Mondays, and 5 Tuesdays. Which day of the week does the 1st of
that month fall on, and how many days does the month have?''
\end{tipbox}
\textbf{Method:} A weekday occurs 5 times in a month only if the month has at least 29 days \emph{and}
the month's first few days include that weekday enough times. For three consecutive weekdays (Sun,
Mon, Tue) to each occur 5 times, the month must have 31 days, and the 1st must fall on the
\emph{first} of the three repeated weekdays. \textbf{The 1st is a Sunday}, and the month has
\textbf{31 days}. (See Trick 10 below for why ``days $-$ 28'' is the fast way to check this.)

\subsection*{I8. Missing Calendar (Given One Fact, Reconstruct the Month)}
\begin{tipbox}[Phrasing]
``In a certain month, 1 January is a Friday. What day of the week is 1 February in a non-leap year?''
\end{tipbox}
\textbf{Method:} January has 31 days $\Rightarrow 31 \bmod 7=3$ odd days $\Rightarrow$ Friday + 3 =
\textbf{Monday}. This is the same relative-arithmetic idea as
\texttt{02-fast-tricks-and-shortcuts.tex}, just applied to ``reconstruct the next month's start''
instead of a specific date within the same month.

\subsection*{I9. Calendar Table / Grid Interpretation}
\begin{tipbox}[Phrasing]
``You are shown a calendar grid for a given month where 1st falls on a Wednesday. What date is the
third Saturday of that month?''
\end{tipbox}
\textbf{Method:} Same as the Nth-weekday method in \texttt{03-question-pattern-catalog.tex} Category
C --- find the first occurrence of the target weekday, then add 7 for each subsequent occurrence.
1st is Wednesday $\Rightarrow$ first Saturday is the 4th $\Rightarrow$ third Saturday
$=4+7+7=18$. \textbf{18th}. When the question shows an actual calendar image rather than describing
it, the method is identical --- just read the first occurrence directly off the grid instead of
computing it.

\subsection*{I10. Date Validation}
\begin{tipbox}[Phrasing]
``Which of these dates cannot exist: 31 April, 30 February, 29 February 2023, 31 June?''
\end{tipbox}
\textbf{Method:} Check each date against the fixed days-per-month rule (30 days have September,
April, June, and November) and the leap-year rule for 29 February. \textbf{31 April, 30 February, 29
February 2023, and 31 June are all invalid} --- April, June, September, and November never have a
31st, no month has a 30th of February, and 2023 is not a leap year so it has no 29 February. This is
an easy, almost purely-recall question type --- don't overthink it.

\section*{Extra Trigger-Word Tricks}

\begin{tipbox}[Trick 6 --- Only Weekdays Are Mentioned (No Dates)]
If a question never gives an actual date, just a weekday and a day-count, skip the full Anchor-Day
Method entirely: take the day count $\bmod 7$ and shift directly from the given weekday.
\end{tipbox}

\begin{tipbox}[Trick 7 --- Same Month Two Dates]
If both dates fall in the same month, never recompute the calendar from the 1st. Just take the
difference between the two date numbers and reduce it $\bmod 7$ (used in I2).
\end{tipbox}

\begin{tipbox}[Trick 8 --- Same Date Different Year]
When comparing the same date across two different years, immediately check whether 29 February falls
in between. If it does, the shift is 2 odd days per year crossed that includes a leap day; otherwise
it's 1 odd day per year.
\end{tipbox}

\begin{tipbox}[Trick 9 --- Nth Weekday of a Month]
Never try to multiply directly to the Nth occurrence. Always find the \emph{first} occurrence of the
target weekday, then add 7 for each additional occurrence needed (used in I6, I7, I9).
\end{tipbox}

\begin{tipbox}[Trick 10 --- ``Which Weekdays Occur 5 Times'']
Don't draw out the whole calendar. A month has (days $-$ 28) ``extra'' days beyond four full weeks,
and exactly that many \emph{consecutive} weekdays (starting from the 1st) occur 5 times. A 31-day
month has 3 such weekdays; a 30-day month has 2; a 28-day month (February, non-leap) has none; a
29-day month (leap February) has exactly 1.
\end{tipbox}

\section*{Decision Tree --- Which Method to Use}

\noindent Use this under time pressure to skip straight to the right method instead of considering
multiple approaches:

\begin{formulabox}[Decision Tree]
\begin{itemize}[leftmargin=1.4em]
  \item \textbf{Only weekdays, no dates?} $\to$ N mod 7 (Trick 6)
  \item \textbf{Two dates in the same month?} $\to$ Difference mod 7 (Trick 7)
  \item \textbf{A specific date, asks for the day it falls on?} $\to$ Anchor-Day Method (\S1 of
  Booklet 02)
  \item \textbf{Leap year status affecting a year-to-year shift?} $\to$ Check if 29 Feb lies
  between (Trick 8)
  \item \textbf{A month name / month code directly?} $\to$ Month Code Method (\S2 of Booklet 02)
  \item \textbf{The Nth or last occurrence of a weekday?} $\to$ First occurrence, then add
  multiples of 7 (Trick 9)
  \item \textbf{Whether a weekday occurs 5 times in a month?} $\to$ days $-$ 28 (Trick 10)
  \item \textbf{Whether two years share the same calendar?} $\to$ Running total of odd days
  (Catalog, Category D)
  \item \textbf{An age, birthday, or leap-year count as the unknown?} $\to$ Odd-day accumulation
  equation (I1, I3)
\end{itemize}
\end{formulabox}

\section*{What This File Deliberately Skips}

Consistent with the Time \& Work module's philosophy: this file does not cover exotic multi-calendar
conversions (Islamic/Hindu calendar mapping), astronomical leap-second edge cases, or Gregorian vs.
Julian calendar transition-date puzzles. These essentially never appear in TCS NQT or general
placement-level Calendar questions.

\section*{Updated Priority Order}

With this file added, the full priority order across all four booklets becomes:

\begin{enumerate}[leftmargin=1.4em]
  \item \textbf{A} --- Direct day lookup (base method --- Anchor-Day)
  \item \textbf{B} --- Relative day arithmetic (no calendar needed)
  \item \textbf{I2} --- Date mirror questions (same family as B)
  \item \textbf{C} --- Nth / last weekday of a month
  \item \textbf{I6/I9} --- First/last working day, calendar table interpretation (same family as C)
  \item \textbf{D} --- Same calendar / repeating year
  \item \textbf{E} --- Leap year reasoning
  \item \textbf{I1/I3} --- Age + calendar, consecutive leap-year logic (same family as E)
  \item \textbf{F} --- Date-to-date and span problems
  \item \textbf{I5} --- Number of Sundays/Mondays between two dates (same family as F)
  \item \textbf{G} --- Reverse problems
  \item \textbf{I8} --- Missing calendar / reconstruct the month (same family as G)
  \item \textbf{H} --- Logic puzzles and word problems
  \item \textbf{I7} --- Month identification from weekday counts (same family as H)
  \item \textbf{I4} --- Financial/office calendar working-day count
  \item \textbf{I10} --- Date validation (lowest effort, easy marks --- don't skip it)
\end{enumerate}

\noindent At this point Calendar coverage is effectively complete for TCS NQT and most general
placement tests. The highest-value next step is the same as with Time \& Work: solve 40--60 mixed
Calendar questions under a timer until each phrasing triggers the right branch of the decision tree
instantly, rather than continuing to add more theory.

\end{document}
